\section{Figure-eight curves and simple lifts}
\label{sec:figures}

In this section surfaces are oriented, of finite type, and have a
hyperbolic metric of curvature $-1$, with cusps or geodesic boundary and
no cone singularities. Curves are unoriented free-homotopy classes unless
an orientation is explicitly chosen for a word calculation.

\subsection{The pants trace relation}

Let $\gamma$ be a closed geodesic with exactly one transverse self-intersection,
not a proper iterate. Its two-lobed regular neighborhood is a pair of
pants: the two ends of each lobe are adjacent in the cyclic order at the
crossing. Its cuffs are essential or peripheral; a compressing disk would
give a geodesic monogon or bigon, or make the curve an iterate.
After tightening the cuffs, allow coincident cuff images in the ambient
surface; this occurs on a one-holed torus. If the three cuff lengths are $a,b,c$,
write
\[
 t_a=2\cosh(a/2),\quad t_b=2\cosh(b/2),\quad t_c=2\cosh(c/2).
\]
A cusp has positive geometric trace $2$, regardless of the sign of a
chosen $\mathrm{SL}_2$ lift. For the figure-eight that wraps the $a$
and $b$ cuffs, the trace identity is
\begin{equation}
 \boxed{t_\gamma=t_at_b+t_c.}
 \label{eq:figuretrace}
\end{equation}
This is Rivin's pants relation~\cite[Theorem 1.3]{RivinFigures}. It also
follows directly from
$\tr(AB^{-1})=\tr A\,\tr B-\tr(AB)$ with the consistent signs of pants
holonomies: choose $\tr A=t_a$, $\tr B=t_b$, and $\tr(AB)=-t_c$.
Rewriting in positive geometric traces produces the plus sign in
\eqref{eq:figuretrace}. One cannot make this replacement in an arbitrary
character identity; Section~\ref{sec:coefficientdistributions} retains
signed traces and therefore has a different sign.

For pants with lengths $(a,b,0)$, the two figures that use the cusp as one
lobe have traces
\begin{equation}
 u=2t_a+t_b,\qquad v=t_a+2t_b,
 \qquad
 t_a=\frac{2u-v}{3},\quad t_b=\frac{2v-u}{3}.
 \label{eq:twofigures}
\end{equation}
Consequently any weight assigned to this pair of cuffs can be rewritten
as a weight in the two associated figure-eight lengths. This gives a
direct reindexing of a pants identity. It produces a paired-curve weight;
it does not automatically produce a universal function of one figure-eight
length.

\subsection{What the local double cover accomplishes}

Use pants generators $a,b$, with third cuff $c=(ab)^{-1}$ and figure
$\gamma=ab^{-1}$. Consider the connected double cover given by
\[
 \chi(a)=\chi(b)=1\in\ZZ/2\ZZ.
\]
Then $\chi(c)=\chi(\gamma)=0$. The $a,b$ cuffs each lift with degree two,
while $c$ has two degree-one lifts. The upstairs surface is a four-holed
sphere with cuff lengths
\begin{equation}
 (2a,2b,c,c).
 \label{eq:doublepants}
\end{equation}
The figure-eight has two closed degree-one lifts, each of which is simple.
One way to see this is to lift its two-loop spine: the lifted graph has
two vertices, and either lift uses one $a$ edge and one $b$ edge to form
a simple cycle.

The two lifted cycles meet at both vertices of that graph. Geometrically,
the two simple lifts intersect each other twice, once over each lift of
the original crossing. Thus two components of the lifted domain do not
mean a disjoint multicurve upstairs.

There is a further global issue: this particular character of the pants
group need not extend to the whole surface. A lobe that is zero in
$H_1(S;\ZZ/2\ZZ)$---for example, a separating lobe on a closed surface---
cannot be assigned value one by an ambient character.
Rivin proves a global substitute: a one-double-point geodesic has a simple
lift in a fourfold cover~\cite[Theorem 4.2]{RivinFigures}.

Even when a suitable cover exists, applying an entire simple-curve identity
upstairs and projecting downstairs introduces other projected topological
types, covering degrees, and multiplicities. Desingularization supplies
an important geometric tool, but selecting just the desired downstairs
curves requires additional accounting.

\subsection{The two types on a punctured torus}

On a once-punctured torus, a figure-eight has a unique essential simple
curve $\alpha$ disjoint from it. To see the topological restriction,
truncate the cusp and consider the complement of its incompressible
pants neighborhood in the resulting compact torus. That complement has
Euler characteristic zero and no disk components, so it consists of two
annuli. One joins a pants cuff to the cusp boundary; the other joins the
remaining two cuffs, which are therefore parallel essential curves.
Conversely, any essential simple
curve disjoint from the figure can be pushed out of its regular
neighborhood and is parallel to these cuffs. Thus its pants neighborhood
has one peripheral cuff and two essential cuffs that tighten to $\alpha$.
Cutting along $\alpha$ gives pants of lengths $(a,a,0)$, where
$a=\ell_\alpha$ and $t=2\cosh(a/2)>2$.

There are exactly three unoriented figure-eight classes in these cut pants.
Indeed, splitting at the unique crossing gives two simple based lobes;
their two distinct peripheral classes determine a choice of two of the three
pants cuffs. Traversing them with opposite boundary orientations is the
one-crossing class; traversing their product with the other relative
orientation is the third cuff. This gives the three choices below:
\begin{center}
\begin{tabular}{@{}p{0.42\textwidth}c c@{}}
\toprule
Lobes & Number & Positive trace\\
\midrule
One essential cuff and the cusp & $2$ & $3t$\\
Both essential cuff copies & $1$ & $t^2+2$\\
\bottomrule
\end{tabular}
\end{center}
The first two are exchanged by the hyperelliptic involution; call their
family $\mathcal N$. The third is hyperelliptic; call its family
$\mathcal H$. Equal lengths in the first row do not identify the two
free-homotopy classes. In torus homology the first family has a primitive
class up to sign and the second twice a primitive class, although its
geodesics are not proper iterates.

For a word check, glue pants generators by $b=s a^{-1}s^{-1}$. The first
two figures can be represented by $a^2b$ and $ab^2$, while the third is
$ab^{-1}=asa s^{-1}$. The last word is not a square. The common disjoint
class $\alpha$ is recoverable from the figure's pants neighborhood, so
grouping all figure-eights by $\alpha$ has no duplication.
More explicitly, the first two words become $a^2sa^{-1}s^{-1}$ and
$asa^{-2}s^{-1}$ in the free basis $a,s$. Their cyclic words are neither
conjugate nor inverse-conjugate. Their homology classes are $[a]$ and
$-[a]$; the third has class $2[a]$ and cyclic word $asa s^{-1}$, which
is not a power of a shorter cyclic word. This also checks primitivity.

Lei--Zhang correct the index set in Rivin's earlier torus identity: its
stated formula accounts for $\mathcal N$, while $\mathcal H$ requires
separate treatment~\cite[p.~2350049-1, abstract]{LeiZhang}. The
classification and weights used here have been derived above; this
citation does not substitute for a check of the full correction paper.
The separation is important here
because the next construction is intended to include \emph{both} families.

\subsection{A positive single-length kernel for all figure-eights}

Define, for $t>2$,
\begin{equation}
 h(t)=1-\sqrt{1-\frac4{t^2}}.
 \label{eq:htrace}
\end{equation}
If $t=2\cosh(\ell/2)$, then $h(t)=2/(1+e^\ell)$. Hence McShane's
identity is $\sum_\alpha h(t_\alpha)=1$. The classification immediately
gives the two separate identities
\begin{align}
 \sum_{\gamma\in\mathcal N}
       \left(1-\sqrt{1-\frac{36}{t_\gamma^2}}\right)&=2,
       \label{eq:nidentity}\\
 \sum_{\gamma\in\mathcal H}
       \left(1-\sqrt{1-\frac4{t_\gamma-2}}\right)&=1.
       \label{eq:hidentity}
\end{align}

To use the same kernel on every one-crossing geodesic, let
\begin{equation}
 \beta(x)=\frac{x^2}{9}+2,\qquad
 F(x)=\sum_{n=0}^{\infty}(-1)^n2^{-n-1}
             h\left(\frac{\beta^{\circ n}(x)}3\right),\qquad x>6.
 \label{eq:universalkernel}
\end{equation}
Here $\beta^{\circ n}$ means iteration, with $\beta^{\circ0}(x)=x$.

\begin{proposition}
On a complete finite-area once-punctured hyperbolic torus,
\begin{equation}
 \boxed{\sum_{\substack{\gamma\text{ unoriented, not iterated}\\
                       i(\gamma,\gamma)=1}}
        F\bigl(2\cosh(\ell_\gamma/2)\bigr)=1.}
 \label{eq:allfigureidentity}
\end{equation}
The kernel is positive and satisfies
\begin{equation}
 \frac14h(x/3)\le F(x)\le\frac12h(x/3),
 \qquad F(x)\sim\frac9{x^2}\quad(x\longrightarrow\infty).
 \label{eq:fbound}
\end{equation}
\end{proposition}

\begin{proof}
For $x>6$,
$\beta(x)-x=(x-3)(x-6)/9>0$. Since $h$ is positive and decreasing,
the absolute values of the summands in~\eqref{eq:universalkernel} decrease,
and the series is absolutely convergent, bounded termwise by $2^{-n-1}$.
Its alternating-series bounds place $F$ between the first term and the
first term minus the second. The second is at most half the first, proving
\eqref{eq:fbound}'s inequalities.

Shifting the series index yields
\[
 2F(x)+F(\beta(x))=h(x/3).
\]
For $x=3t$ this reads
\begin{equation}
 2F(3t)+F(t^2+2)=h(t).
 \label{eq:kernelrelation}
\end{equation}
The three figure-eights associated with a simple $\alpha$ have exactly
the traces on the left. Sum over $\alpha$ and use McShane's identity.
Positivity justifies the regrouping and proves convergence of the full
geodesic sum. Finally $h(x/3)=18x^{-2}+O(x^{-4})$ and the second and
subsequent terms contribute $O(x^{-4})$, giving the stated asymptotic.
\end{proof}

The same alternating-series argument supplies an explicit finite
evaluation contract. For $N\geq1$ let $F_N$ be the sum of the first $N$
terms in~\eqref{eq:universalkernel}, and put $x_N=\beta^{\circ N}(x)$.
Then
\begin{equation}
 0\leq(-1)^N\bigl(F(x)-F_N(x)\bigr)
 \leq 2^{-N-1}h(x_N/3)\leq 2^{-N-1}.
 \label{eq:figurekerneltail}
\end{equation}
Moreover the exact finite identity
\begin{equation}
 2F_N(x)+F_N(\beta(x))-h(x/3)
   =(-1)^{N-1}2^{-N}h(x_N/3)
 \label{eq:figurefiniteidentity}
\end{equation}
provides a check independent of extrapolating a computed infinite sum.
For large $t$, evaluate
$h(t)=4/[t^2(1+\sqrt{1-4/t^2})]$ to avoid subtraction of nearby numbers.
These bounds concern evaluation of the kernel. Approximating the sum over
surface curves additionally requires a complete index enumeration and a
bound for the omitted curve classes; no such numerical approximation is
asserted here.

This is an elementary construction from the corrected classification,
not a claim of a new identity in the literature. It is restricted to the
cusped torus. Cone continuation requires checking existence of the indexed
geodesics as well as continuing the weight; the cone obstruction discussed
by Lei--Zhang is a reason to keep that distinction explicit~\cite{LeiZhang}.

\paragraph{Computational support.}
Groups G2--G3 of \path{experiments/quantum_geodesics/README.md}
check finite signed-trace identities, the distinct torus labels,
and the universal kernel's finite coefficients and tails.
The exact telescoping coefficients are checked algebraically;
at large truncation order a cancellation-limited floating residual
is not an independent relative check of a much smaller positive tail.
These fixtures do not enumerate all curves or approximate the full
surface identity. Their result path is
\path{artifacts/quantum_geodesics/result.json}, with claim pointers in
\path{book/claims/sections17-22.json}.
